% ============================================================
%  APPENDIX A — DATASET DOCUMENTATION (DATASHEET)
%
%  Structured following Gebru et al. (2021) \cite{gebru2021datasheets}
%  seven-section framework. Place after \section{Conclusion} in the
%  main .tex file with:
%    \appendix
%    % ============================================================
%  APPENDIX A — DATASET DOCUMENTATION (DATASHEET)
%
%  Structured following Gebru et al. (2021) \cite{gebru2021datasheets}
%  seven-section framework. Place after \section{Conclusion} in the
%  main .tex file with:
%    \appendix
%    % ============================================================
%  APPENDIX A — DATASET DOCUMENTATION (DATASHEET)
%
%  Structured following Gebru et al. (2021) \cite{gebru2021datasheets}
%  seven-section framework. Place after \section{Conclusion} in the
%  main .tex file with:
%    \appendix
%    % ============================================================
%  APPENDIX A — DATASET DOCUMENTATION (DATASHEET)
%
%  Structured following Gebru et al. (2021) \cite{gebru2021datasheets}
%  seven-section framework. Place after \section{Conclusion} in the
%  main .tex file with:
%    \appendix
%    \input{appendix_datasheet}
%
%  New \bibitem key:
%    gebru2021datasheets
%
%  Required preamble additions:
%    \usepackage{enumitem}       % for \begin{enumerate}[label=\textbf{Q\arabic*.}]
%    \usepackage{xcolor}
%    \definecolor{shadecolor}{RGB}{245,245,245}
%    \usepackage{framed}         % for \begin{shaded}...\end{shaded} question boxes
%    (or replace shaded blocks with \noindent\textbf{...} if framed not available)
% ============================================================

\section{Dataset Documentation}
\label{sec:datasheet}

This appendix provides a structured documentation record for
\textbf{\datasetname} following the \emph{Datasheets for Datasets}
framework of Gebru~\etal~\cite{gebru2021datasheets}.
We address all seven canonical categories: motivation, composition,
collection process, preprocessing and labelling, intended uses,
distribution, and maintenance.
This documentation is also provided in machine-readable form as
\texttt{datasheet.json} in the dataset repository.

%-------------------------------------------------------------------
\subsection*{A.1 \quad Motivation}
\label{ssec:ds_motivation}
%-------------------------------------------------------------------

\paragraph{Purpose.}
\textbf{\datasetname} was created to support the development and
evaluation of AI-generated image detectors under conditions that
eliminate the two principal dataset biases identified in prior work:
JPEG compression asymmetry between real and AI images, and semantic
content mismatch between the two classes.
No comparable leak-free paired dataset existed at the time of
construction that covered more than four generators or more than one
backbone architecture family.

\paragraph{Creators.}
The dataset was constructed by Shanmukesh Bonala
(\texttt{Shanmuk4622}) as part of doctoral research at VIT-AP
University, under the supervision of Dr.~Deepthi Godavarthi.
The build pipeline was implemented across eleven version-controlled
Kaggle notebooks (NB00--NB11) available at the repository URL.

\paragraph{Funding.}
No external funding was received for this work.
Computational resources were provided by Kaggle's free-tier GPU
programme (dual NVIDIA Tesla T4).

%-------------------------------------------------------------------
\subsection*{A.2 \quad Composition}
\label{ssec:ds_composition}
%-------------------------------------------------------------------

\paragraph{What the dataset represents.}
\textbf{\datasetname} consists of 70,000 RGB images at $512 \times 512$
pixels, divided into a real class (10,000 photographs) and an AI
class (60,000 synthetic images produced by six text-to-image diffusion
models).
Each real image has exactly six AI partners, one per generator,
all generated from the same BLIP-2 caption derived from that real
image.

\paragraph{Instances.}
Each instance is a single image stored as a lossless PNG in an Apache
Parquet shard.
Metadata fields (listed in full in Table~\ref{tab:schema}) include a
globally unique image identifier, the pairing key
(\texttt{source\_real\_id}), label (0~=~real, 1~=~AI), generator key,
source dataset name, generation prompt, spatial dimensions, original
pre-crop dimensions, generation hyperparameters, caption model name,
pipeline version, SHA-256 digest of the PNG bytes, and a UTC creation
timestamp.

\paragraph{Split counts.}
The dataset is partitioned into a training split (49,392 images,
7,056 pairs), a validation split (10,122 images, 1,446 pairs), and a
test split (10,486 images, 1,498 pairs).
Splits are assigned deterministically at the pair level
(Section~\ref{ssec:splits}) so that no real image and its six AI
partners appear in different folds.

\paragraph{Label noise.}
All real images carry label~0 and all AI images carry label~1 by
construction.
No human annotation was involved in label assignment; labels are
determined entirely by the data source.
Label noise is therefore zero by design, but generation failures
(e.g., a model collapsing to a degenerate output) may produce
images that are visually uninformative.
One such group was identified during deduplication (four W\"{u}rstchen
images sharing identical byte content, flagged in the manifest via
shared SHA-256 digest) and is noted in Section~\ref{ssec:dedup}.

\paragraph{Sensitive content.}
Real images are sourced from MS-COCO and ImageNet validation splits.
COCO images are Flickr-sourced photographs under Creative Commons
licences; they may contain depictions of people, including children,
in public settings.
AI images are synthetic outputs and do not depict real individuals.
The dataset contains no sexually explicit content; generation
pipelines were run with \texttt{safety\_checker=None} solely to
prevent the safety filter from introducing a class-asymmetric
post-processing step, not to produce unsafe content.

\paragraph{Self-contained.}
Image bytes are embedded directly in the Parquet files; no external
URLs or file paths are required to reproduce the dataset.

%-------------------------------------------------------------------
\subsection*{A.3 \quad Collection Process}
\label{ssec:ds_collection}
%-------------------------------------------------------------------

\paragraph{Real images.}
5,000 images were sampled from the MS-COCO 2017 validation
split~\cite{lin2014coco} and 5,000 from the ImageNet-1k validation
split~\cite{deng2009imagenet}, downloaded via the canonical
HuggingFace dataset repositories.
No web scraping was performed; all real images were obtained from
their official public releases.
Sampling was deterministic: images were drawn in index order up to
the target count, with MASTER\_SEED~$= 42$ controlling any shuffled
ordering.

\paragraph{AI images.}
Synthetic images were generated programmatically during June 2026
on Kaggle GPU sessions (NVIDIA Tesla T4, 15.6~GB VRAM) using the
HuggingFace Diffusers library.
Six generation notebooks (NB03--NB08) ran in parallel across multiple
Kaggle accounts using a hash-based worker-partitioning scheme to
avoid shard collisions.
No human selection was applied to generated images; every output
produced by the pipeline was included, subject only to the
integrity checks described in NB10.

\paragraph{Captions.}
Captions were generated in NB02 by feeding each real image (after
canonical preprocessing) to BLIP-2 (\texttt{Salesforce/blip2-opt-2.7b})
running beam search (5 beams, max 60 new tokens, repetition penalty
1.3), then cleaning and token-capping the result as described in
Section~\ref{ssec:captions}.
No human captioning was performed.

\paragraph{Workers and consent.}
All data collection and generation was automated; no human workers
or crowdworkers were involved.

%-------------------------------------------------------------------
\subsection*{A.4 \quad Preprocessing, Cleaning, and Labelling}
\label{ssec:ds_preprocessing}
%-------------------------------------------------------------------

\paragraph{Preprocessing.}
Every image — real or AI — was passed through the canonical
preprocessing pipeline (Algorithm~\ref{alg:preprocess}) with no
branching on the class label.
The pipeline performs EXIF transpose, RGB conversion, centre-crop to
square, Lanczos resize to $512 \times 512$, JPEG equalisation at
quality factor 95 with 4:4:4 chroma subsampling, and lossless PNG
serialisation at compress level 6.

\paragraph{Quality checks.}
Every stored image was verified against the \texttt{png\_is\_canonical}
predicate (PNG format, RGB mode, $512 \times 512$ dimensions, no ICC
profile) and the stored SHA-256 digest.
Images failing these checks were excluded and the corresponding
generation run was re-executed.
Zero images in the released dataset fail either check.

\paragraph{Labelling.}
Labels are assigned automatically by provenance: images sourced from
COCO or ImageNet receive label~0; images output by any of the six
generators receive label~1.
No human labelling was performed.

\paragraph{Raw data availability.}
The raw (pre-preprocessed) real images are available from their
respective original dataset repositories.
Raw generator outputs prior to canonical preprocessing are not
retained in the repository; only the preprocessed canonical form
is stored.
The preprocessing code is fully open-source and deterministic, so
raw generator outputs can be reproduced from the stored prompts
and seeds.

%-------------------------------------------------------------------
\subsection*{A.5 \quad Uses}
\label{ssec:ds_uses}
%-------------------------------------------------------------------

\paragraph{Intended uses.}
\textbf{\datasetname} is intended for:
\begin{itemize}[leftmargin=1.5em, itemsep=2pt]
  \item Training and evaluating binary real-vs-AI image classifiers
    under bias-free conditions.
  \item Studying per-generator detectability gradients and
    cross-generator generalisation.
  \item Benchmarking novel detection architectures on a dataset where
    JPEG and resolution shortcuts are provably eliminated.
  \item Ablation studies on the effect of preprocessing normalisation
    on detector performance.
\end{itemize}

\paragraph{Uses to avoid.}
\begin{itemize}[leftmargin=1.5em, itemsep=2pt]
  \item Training detectors intended for deployment on images that have
    not been similarly preprocessed, without re-evaluating robustness
    to the omitted preprocessing.
  \item Using the \texttt{orig\_width} and \texttt{orig\_height} fields
    as training features; these encode source-level provenance and
    reintroduce a resolution-based shortcut.
  \item Any application requiring a larger or more demographically
    representative real-image corpus than COCO and ImageNet provide.
  \item Commercial deployment of detectors trained solely on this
    dataset without further validation, given the non-commercial
    licence restriction on the ImageNet component.
\end{itemize}

\paragraph{Impact of misuse.}
A detector trained under conditions that inflate performance
(e.g., on an uncorrected dataset) may provide false assurance
of reliability when deployed on images that do not share those
conditions.
We caution that \textbf{\datasetname} itself is a benchmark, not a
production detection system, and that all reported metrics should
be interpreted within the experimental conditions described in
Section~\ref{sec:experiments}.

%-------------------------------------------------------------------
\subsection*{A.6 \quad Distribution}
\label{ssec:ds_distribution}
%-------------------------------------------------------------------

\paragraph{Access.}
The dataset is publicly available at:\\
\texttt{https://huggingface.co/datasets/Shanmuk4622/ai-detection-dataset-v2}\\
All images, metadata, manifest, splits, validation report, and
generation configuration are accessible without authentication.

\paragraph{Licence.}
The dataset inherits the most restrictive licence applicable to any
of its components:
\begin{itemize}[leftmargin=1.5em, itemsep=2pt]
  \item \textbf{COCO images}: Flickr photographs under Creative Commons
    licences (primarily CC BY 2.0). Commercial use permitted.
  \item \textbf{ImageNet images}: ILSVRC terms of use.
    \emph{Non-commercial research use only.}
  \item \textbf{AI images}: Synthetic outputs of the six generators,
    each under its own open model licence. All six permit
    non-commercial research use.
\end{itemize}
The effective licence for the combined dataset is therefore
\textbf{non-commercial research use only}.
This is not legal advice; users are responsible for verifying
compliance with the licence terms of each component.

\paragraph{DOI and versioning.}
The dataset is tagged at version \texttt{v2.0} in the HuggingFace
repository.
A DOI will be assigned upon journal acceptance through the
HuggingFace dataset DOI programme.

\paragraph{Export controls.}
No export control or other legal restrictions beyond the licence
terms above are known to apply to this dataset.

%-------------------------------------------------------------------
\subsection*{A.7 \quad Maintenance}
\label{ssec:ds_maintenance}
%-------------------------------------------------------------------

\paragraph{Maintainer.}
The dataset is maintained by the first author,
\texttt{Shanmuk4622} on HuggingFace.
Contact: via the HuggingFace repository discussion tab.

\paragraph{Erratum policy.}
Errors in image data, metadata, or splits discovered after release
will be documented in the repository's \texttt{ERRATA.md} file and
will not silently overwrite existing files.
Corrected versions will be tagged with an incremented minor version
number (e.g., \texttt{v2.1}) with a changelog entry.

\paragraph{Updates.}
No scheduled updates are planned.
Extensions (additional generators, larger image counts, or
higher resolution variants) will be released as separate versioned
datasets rather than by modifying the released \texttt{v2.0} tag,
preserving the reproducibility of published results.

\paragraph{Reproducibility.}
The complete build is reproducible from:
\begin{enumerate}[leftmargin=1.5em, itemsep=2pt]
  \item The eleven generation notebooks (NB00--NB11) stored in the
    repository under \texttt{notebooks/}.
  \item The shared library \texttt{ai\_detect\_common.py}
    (pipeline version 1.2).
  \item The configuration file \texttt{config.json}, which records
    all hyperparameters including \texttt{MASTER\_SEED}~$= 42$.
\end{enumerate}
Changing \texttt{MASTER\_SEED} in NB00 and re-running the full
pipeline produces a statistically independent dataset with the same
structure, enabling sensitivity analysis over the random seed.

\paragraph{Legacy support.}
\texttt{v1.x} of the dataset (50,000 images, different repository)
remains available at its original HuggingFace URL for continuity
with any work that cited it prior to this paper.
It is not recommended for new work; \texttt{v2.0} supersedes it
in all respects.

% ─── New bibliography entry ─────────────────────────────────────
% Add to thebibliography block in related_work_and_references.tex

%\bibitem{gebru2021datasheets}
%T.~Gebru, J.~Morgenstern, B.~Vecchione, J.~W.~Vaughan, H.~Wallach,
%H.~Daum\'{e}~III, and K.~Crawford,
%``Datasheets for datasets,''
%\emph{Communications of the ACM}, vol.~64, no.~12, pp.~86--92, 2021.
%doi:10.1145/3458723.

%
%  New \bibitem key:
%    gebru2021datasheets
%
%  Required preamble additions:
%    \usepackage{enumitem}       % for \begin{enumerate}[label=\textbf{Q\arabic*.}]
%    \usepackage{xcolor}
%    \definecolor{shadecolor}{RGB}{245,245,245}
%    \usepackage{framed}         % for \begin{shaded}...\end{shaded} question boxes
%    (or replace shaded blocks with \noindent\textbf{...} if framed not available)
% ============================================================

\section{Dataset Documentation}
\label{sec:datasheet}

This appendix provides a structured documentation record for
\textbf{\datasetname} following the \emph{Datasheets for Datasets}
framework of Gebru~\etal~\cite{gebru2021datasheets}.
We address all seven canonical categories: motivation, composition,
collection process, preprocessing and labelling, intended uses,
distribution, and maintenance.
This documentation is also provided in machine-readable form as
\texttt{datasheet.json} in the dataset repository.

%-------------------------------------------------------------------
\subsection*{A.1 \quad Motivation}
\label{ssec:ds_motivation}
%-------------------------------------------------------------------

\paragraph{Purpose.}
\textbf{\datasetname} was created to support the development and
evaluation of AI-generated image detectors under conditions that
eliminate the two principal dataset biases identified in prior work:
JPEG compression asymmetry between real and AI images, and semantic
content mismatch between the two classes.
No comparable leak-free paired dataset existed at the time of
construction that covered more than four generators or more than one
backbone architecture family.

\paragraph{Creators.}
The dataset was constructed by Shanmukesh Bonala
(\texttt{Shanmuk4622}) as part of doctoral research at VIT-AP
University, under the supervision of Dr.~Deepthi Godavarthi.
The build pipeline was implemented across eleven version-controlled
Kaggle notebooks (NB00--NB11) available at the repository URL.

\paragraph{Funding.}
No external funding was received for this work.
Computational resources were provided by Kaggle's free-tier GPU
programme (dual NVIDIA Tesla T4).

%-------------------------------------------------------------------
\subsection*{A.2 \quad Composition}
\label{ssec:ds_composition}
%-------------------------------------------------------------------

\paragraph{What the dataset represents.}
\textbf{\datasetname} consists of 70,000 RGB images at $512 \times 512$
pixels, divided into a real class (10,000 photographs) and an AI
class (60,000 synthetic images produced by six text-to-image diffusion
models).
Each real image has exactly six AI partners, one per generator,
all generated from the same BLIP-2 caption derived from that real
image.

\paragraph{Instances.}
Each instance is a single image stored as a lossless PNG in an Apache
Parquet shard.
Metadata fields (listed in full in Table~\ref{tab:schema}) include a
globally unique image identifier, the pairing key
(\texttt{source\_real\_id}), label (0~=~real, 1~=~AI), generator key,
source dataset name, generation prompt, spatial dimensions, original
pre-crop dimensions, generation hyperparameters, caption model name,
pipeline version, SHA-256 digest of the PNG bytes, and a UTC creation
timestamp.

\paragraph{Split counts.}
The dataset is partitioned into a training split (49,392 images,
7,056 pairs), a validation split (10,122 images, 1,446 pairs), and a
test split (10,486 images, 1,498 pairs).
Splits are assigned deterministically at the pair level
(Section~\ref{ssec:splits}) so that no real image and its six AI
partners appear in different folds.

\paragraph{Label noise.}
All real images carry label~0 and all AI images carry label~1 by
construction.
No human annotation was involved in label assignment; labels are
determined entirely by the data source.
Label noise is therefore zero by design, but generation failures
(e.g., a model collapsing to a degenerate output) may produce
images that are visually uninformative.
One such group was identified during deduplication (four W\"{u}rstchen
images sharing identical byte content, flagged in the manifest via
shared SHA-256 digest) and is noted in Section~\ref{ssec:dedup}.

\paragraph{Sensitive content.}
Real images are sourced from MS-COCO and ImageNet validation splits.
COCO images are Flickr-sourced photographs under Creative Commons
licences; they may contain depictions of people, including children,
in public settings.
AI images are synthetic outputs and do not depict real individuals.
The dataset contains no sexually explicit content; generation
pipelines were run with \texttt{safety\_checker=None} solely to
prevent the safety filter from introducing a class-asymmetric
post-processing step, not to produce unsafe content.

\paragraph{Self-contained.}
Image bytes are embedded directly in the Parquet files; no external
URLs or file paths are required to reproduce the dataset.

%-------------------------------------------------------------------
\subsection*{A.3 \quad Collection Process}
\label{ssec:ds_collection}
%-------------------------------------------------------------------

\paragraph{Real images.}
5,000 images were sampled from the MS-COCO 2017 validation
split~\cite{lin2014coco} and 5,000 from the ImageNet-1k validation
split~\cite{deng2009imagenet}, downloaded via the canonical
HuggingFace dataset repositories.
No web scraping was performed; all real images were obtained from
their official public releases.
Sampling was deterministic: images were drawn in index order up to
the target count, with MASTER\_SEED~$= 42$ controlling any shuffled
ordering.

\paragraph{AI images.}
Synthetic images were generated programmatically during June 2026
on Kaggle GPU sessions (NVIDIA Tesla T4, 15.6~GB VRAM) using the
HuggingFace Diffusers library.
Six generation notebooks (NB03--NB08) ran in parallel across multiple
Kaggle accounts using a hash-based worker-partitioning scheme to
avoid shard collisions.
No human selection was applied to generated images; every output
produced by the pipeline was included, subject only to the
integrity checks described in NB10.

\paragraph{Captions.}
Captions were generated in NB02 by feeding each real image (after
canonical preprocessing) to BLIP-2 (\texttt{Salesforce/blip2-opt-2.7b})
running beam search (5 beams, max 60 new tokens, repetition penalty
1.3), then cleaning and token-capping the result as described in
Section~\ref{ssec:captions}.
No human captioning was performed.

\paragraph{Workers and consent.}
All data collection and generation was automated; no human workers
or crowdworkers were involved.

%-------------------------------------------------------------------
\subsection*{A.4 \quad Preprocessing, Cleaning, and Labelling}
\label{ssec:ds_preprocessing}
%-------------------------------------------------------------------

\paragraph{Preprocessing.}
Every image — real or AI — was passed through the canonical
preprocessing pipeline (Algorithm~\ref{alg:preprocess}) with no
branching on the class label.
The pipeline performs EXIF transpose, RGB conversion, centre-crop to
square, Lanczos resize to $512 \times 512$, JPEG equalisation at
quality factor 95 with 4:4:4 chroma subsampling, and lossless PNG
serialisation at compress level 6.

\paragraph{Quality checks.}
Every stored image was verified against the \texttt{png\_is\_canonical}
predicate (PNG format, RGB mode, $512 \times 512$ dimensions, no ICC
profile) and the stored SHA-256 digest.
Images failing these checks were excluded and the corresponding
generation run was re-executed.
Zero images in the released dataset fail either check.

\paragraph{Labelling.}
Labels are assigned automatically by provenance: images sourced from
COCO or ImageNet receive label~0; images output by any of the six
generators receive label~1.
No human labelling was performed.

\paragraph{Raw data availability.}
The raw (pre-preprocessed) real images are available from their
respective original dataset repositories.
Raw generator outputs prior to canonical preprocessing are not
retained in the repository; only the preprocessed canonical form
is stored.
The preprocessing code is fully open-source and deterministic, so
raw generator outputs can be reproduced from the stored prompts
and seeds.

%-------------------------------------------------------------------
\subsection*{A.5 \quad Uses}
\label{ssec:ds_uses}
%-------------------------------------------------------------------

\paragraph{Intended uses.}
\textbf{\datasetname} is intended for:
\begin{itemize}[leftmargin=1.5em, itemsep=2pt]
  \item Training and evaluating binary real-vs-AI image classifiers
    under bias-free conditions.
  \item Studying per-generator detectability gradients and
    cross-generator generalisation.
  \item Benchmarking novel detection architectures on a dataset where
    JPEG and resolution shortcuts are provably eliminated.
  \item Ablation studies on the effect of preprocessing normalisation
    on detector performance.
\end{itemize}

\paragraph{Uses to avoid.}
\begin{itemize}[leftmargin=1.5em, itemsep=2pt]
  \item Training detectors intended for deployment on images that have
    not been similarly preprocessed, without re-evaluating robustness
    to the omitted preprocessing.
  \item Using the \texttt{orig\_width} and \texttt{orig\_height} fields
    as training features; these encode source-level provenance and
    reintroduce a resolution-based shortcut.
  \item Any application requiring a larger or more demographically
    representative real-image corpus than COCO and ImageNet provide.
  \item Commercial deployment of detectors trained solely on this
    dataset without further validation, given the non-commercial
    licence restriction on the ImageNet component.
\end{itemize}

\paragraph{Impact of misuse.}
A detector trained under conditions that inflate performance
(e.g., on an uncorrected dataset) may provide false assurance
of reliability when deployed on images that do not share those
conditions.
We caution that \textbf{\datasetname} itself is a benchmark, not a
production detection system, and that all reported metrics should
be interpreted within the experimental conditions described in
Section~\ref{sec:experiments}.

%-------------------------------------------------------------------
\subsection*{A.6 \quad Distribution}
\label{ssec:ds_distribution}
%-------------------------------------------------------------------

\paragraph{Access.}
The dataset is publicly available at:\\
\texttt{https://huggingface.co/datasets/Shanmuk4622/ai-detection-dataset-v2}\\
All images, metadata, manifest, splits, validation report, and
generation configuration are accessible without authentication.

\paragraph{Licence.}
The dataset inherits the most restrictive licence applicable to any
of its components:
\begin{itemize}[leftmargin=1.5em, itemsep=2pt]
  \item \textbf{COCO images}: Flickr photographs under Creative Commons
    licences (primarily CC BY 2.0). Commercial use permitted.
  \item \textbf{ImageNet images}: ILSVRC terms of use.
    \emph{Non-commercial research use only.}
  \item \textbf{AI images}: Synthetic outputs of the six generators,
    each under its own open model licence. All six permit
    non-commercial research use.
\end{itemize}
The effective licence for the combined dataset is therefore
\textbf{non-commercial research use only}.
This is not legal advice; users are responsible for verifying
compliance with the licence terms of each component.

\paragraph{DOI and versioning.}
The dataset is tagged at version \texttt{v2.0} in the HuggingFace
repository.
A DOI will be assigned upon journal acceptance through the
HuggingFace dataset DOI programme.

\paragraph{Export controls.}
No export control or other legal restrictions beyond the licence
terms above are known to apply to this dataset.

%-------------------------------------------------------------------
\subsection*{A.7 \quad Maintenance}
\label{ssec:ds_maintenance}
%-------------------------------------------------------------------

\paragraph{Maintainer.}
The dataset is maintained by the first author,
\texttt{Shanmuk4622} on HuggingFace.
Contact: via the HuggingFace repository discussion tab.

\paragraph{Erratum policy.}
Errors in image data, metadata, or splits discovered after release
will be documented in the repository's \texttt{ERRATA.md} file and
will not silently overwrite existing files.
Corrected versions will be tagged with an incremented minor version
number (e.g., \texttt{v2.1}) with a changelog entry.

\paragraph{Updates.}
No scheduled updates are planned.
Extensions (additional generators, larger image counts, or
higher resolution variants) will be released as separate versioned
datasets rather than by modifying the released \texttt{v2.0} tag,
preserving the reproducibility of published results.

\paragraph{Reproducibility.}
The complete build is reproducible from:
\begin{enumerate}[leftmargin=1.5em, itemsep=2pt]
  \item The eleven generation notebooks (NB00--NB11) stored in the
    repository under \texttt{notebooks/}.
  \item The shared library \texttt{ai\_detect\_common.py}
    (pipeline version 1.2).
  \item The configuration file \texttt{config.json}, which records
    all hyperparameters including \texttt{MASTER\_SEED}~$= 42$.
\end{enumerate}
Changing \texttt{MASTER\_SEED} in NB00 and re-running the full
pipeline produces a statistically independent dataset with the same
structure, enabling sensitivity analysis over the random seed.

\paragraph{Legacy support.}
\texttt{v1.x} of the dataset (50,000 images, different repository)
remains available at its original HuggingFace URL for continuity
with any work that cited it prior to this paper.
It is not recommended for new work; \texttt{v2.0} supersedes it
in all respects.

% ─── New bibliography entry ─────────────────────────────────────
% Add to thebibliography block in related_work_and_references.tex

%\bibitem{gebru2021datasheets}
%T.~Gebru, J.~Morgenstern, B.~Vecchione, J.~W.~Vaughan, H.~Wallach,
%H.~Daum\'{e}~III, and K.~Crawford,
%``Datasheets for datasets,''
%\emph{Communications of the ACM}, vol.~64, no.~12, pp.~86--92, 2021.
%doi:10.1145/3458723.

%
%  New \bibitem key:
%    gebru2021datasheets
%
%  Required preamble additions:
%    \usepackage{enumitem}       % for \begin{enumerate}[label=\textbf{Q\arabic*.}]
%    \usepackage{xcolor}
%    \definecolor{shadecolor}{RGB}{245,245,245}
%    \usepackage{framed}         % for \begin{shaded}...\end{shaded} question boxes
%    (or replace shaded blocks with \noindent\textbf{...} if framed not available)
% ============================================================

\section{Dataset Documentation}
\label{sec:datasheet}

This appendix provides a structured documentation record for
\textbf{\datasetname} following the \emph{Datasheets for Datasets}
framework of Gebru~\etal~\cite{gebru2021datasheets}.
We address all seven canonical categories: motivation, composition,
collection process, preprocessing and labelling, intended uses,
distribution, and maintenance.
This documentation is also provided in machine-readable form as
\texttt{datasheet.json} in the dataset repository.

%-------------------------------------------------------------------
\subsection*{A.1 \quad Motivation}
\label{ssec:ds_motivation}
%-------------------------------------------------------------------

\paragraph{Purpose.}
\textbf{\datasetname} was created to support the development and
evaluation of AI-generated image detectors under conditions that
eliminate the two principal dataset biases identified in prior work:
JPEG compression asymmetry between real and AI images, and semantic
content mismatch between the two classes.
No comparable leak-free paired dataset existed at the time of
construction that covered more than four generators or more than one
backbone architecture family.

\paragraph{Creators.}
The dataset was constructed by Shanmukesh Bonala
(\texttt{Shanmuk4622}) as part of doctoral research at VIT-AP
University, under the supervision of Dr.~Deepthi Godavarthi.
The build pipeline was implemented across eleven version-controlled
Kaggle notebooks (NB00--NB11) available at the repository URL.

\paragraph{Funding.}
No external funding was received for this work.
Computational resources were provided by Kaggle's free-tier GPU
programme (dual NVIDIA Tesla T4).

%-------------------------------------------------------------------
\subsection*{A.2 \quad Composition}
\label{ssec:ds_composition}
%-------------------------------------------------------------------

\paragraph{What the dataset represents.}
\textbf{\datasetname} consists of 70,000 RGB images at $512 \times 512$
pixels, divided into a real class (10,000 photographs) and an AI
class (60,000 synthetic images produced by six text-to-image diffusion
models).
Each real image has exactly six AI partners, one per generator,
all generated from the same BLIP-2 caption derived from that real
image.

\paragraph{Instances.}
Each instance is a single image stored as a lossless PNG in an Apache
Parquet shard.
Metadata fields (listed in full in Table~\ref{tab:schema}) include a
globally unique image identifier, the pairing key
(\texttt{source\_real\_id}), label (0~=~real, 1~=~AI), generator key,
source dataset name, generation prompt, spatial dimensions, original
pre-crop dimensions, generation hyperparameters, caption model name,
pipeline version, SHA-256 digest of the PNG bytes, and a UTC creation
timestamp.

\paragraph{Split counts.}
The dataset is partitioned into a training split (49,392 images,
7,056 pairs), a validation split (10,122 images, 1,446 pairs), and a
test split (10,486 images, 1,498 pairs).
Splits are assigned deterministically at the pair level
(Section~\ref{ssec:splits}) so that no real image and its six AI
partners appear in different folds.

\paragraph{Label noise.}
All real images carry label~0 and all AI images carry label~1 by
construction.
No human annotation was involved in label assignment; labels are
determined entirely by the data source.
Label noise is therefore zero by design, but generation failures
(e.g., a model collapsing to a degenerate output) may produce
images that are visually uninformative.
One such group was identified during deduplication (four W\"{u}rstchen
images sharing identical byte content, flagged in the manifest via
shared SHA-256 digest) and is noted in Section~\ref{ssec:dedup}.

\paragraph{Sensitive content.}
Real images are sourced from MS-COCO and ImageNet validation splits.
COCO images are Flickr-sourced photographs under Creative Commons
licences; they may contain depictions of people, including children,
in public settings.
AI images are synthetic outputs and do not depict real individuals.
The dataset contains no sexually explicit content; generation
pipelines were run with \texttt{safety\_checker=None} solely to
prevent the safety filter from introducing a class-asymmetric
post-processing step, not to produce unsafe content.

\paragraph{Self-contained.}
Image bytes are embedded directly in the Parquet files; no external
URLs or file paths are required to reproduce the dataset.

%-------------------------------------------------------------------
\subsection*{A.3 \quad Collection Process}
\label{ssec:ds_collection}
%-------------------------------------------------------------------

\paragraph{Real images.}
5,000 images were sampled from the MS-COCO 2017 validation
split~\cite{lin2014coco} and 5,000 from the ImageNet-1k validation
split~\cite{deng2009imagenet}, downloaded via the canonical
HuggingFace dataset repositories.
No web scraping was performed; all real images were obtained from
their official public releases.
Sampling was deterministic: images were drawn in index order up to
the target count, with MASTER\_SEED~$= 42$ controlling any shuffled
ordering.

\paragraph{AI images.}
Synthetic images were generated programmatically during June 2026
on Kaggle GPU sessions (NVIDIA Tesla T4, 15.6~GB VRAM) using the
HuggingFace Diffusers library.
Six generation notebooks (NB03--NB08) ran in parallel across multiple
Kaggle accounts using a hash-based worker-partitioning scheme to
avoid shard collisions.
No human selection was applied to generated images; every output
produced by the pipeline was included, subject only to the
integrity checks described in NB10.

\paragraph{Captions.}
Captions were generated in NB02 by feeding each real image (after
canonical preprocessing) to BLIP-2 (\texttt{Salesforce/blip2-opt-2.7b})
running beam search (5 beams, max 60 new tokens, repetition penalty
1.3), then cleaning and token-capping the result as described in
Section~\ref{ssec:captions}.
No human captioning was performed.

\paragraph{Workers and consent.}
All data collection and generation was automated; no human workers
or crowdworkers were involved.

%-------------------------------------------------------------------
\subsection*{A.4 \quad Preprocessing, Cleaning, and Labelling}
\label{ssec:ds_preprocessing}
%-------------------------------------------------------------------

\paragraph{Preprocessing.}
Every image — real or AI — was passed through the canonical
preprocessing pipeline (Algorithm~\ref{alg:preprocess}) with no
branching on the class label.
The pipeline performs EXIF transpose, RGB conversion, centre-crop to
square, Lanczos resize to $512 \times 512$, JPEG equalisation at
quality factor 95 with 4:4:4 chroma subsampling, and lossless PNG
serialisation at compress level 6.

\paragraph{Quality checks.}
Every stored image was verified against the \texttt{png\_is\_canonical}
predicate (PNG format, RGB mode, $512 \times 512$ dimensions, no ICC
profile) and the stored SHA-256 digest.
Images failing these checks were excluded and the corresponding
generation run was re-executed.
Zero images in the released dataset fail either check.

\paragraph{Labelling.}
Labels are assigned automatically by provenance: images sourced from
COCO or ImageNet receive label~0; images output by any of the six
generators receive label~1.
No human labelling was performed.

\paragraph{Raw data availability.}
The raw (pre-preprocessed) real images are available from their
respective original dataset repositories.
Raw generator outputs prior to canonical preprocessing are not
retained in the repository; only the preprocessed canonical form
is stored.
The preprocessing code is fully open-source and deterministic, so
raw generator outputs can be reproduced from the stored prompts
and seeds.

%-------------------------------------------------------------------
\subsection*{A.5 \quad Uses}
\label{ssec:ds_uses}
%-------------------------------------------------------------------

\paragraph{Intended uses.}
\textbf{\datasetname} is intended for:
\begin{itemize}[leftmargin=1.5em, itemsep=2pt]
  \item Training and evaluating binary real-vs-AI image classifiers
    under bias-free conditions.
  \item Studying per-generator detectability gradients and
    cross-generator generalisation.
  \item Benchmarking novel detection architectures on a dataset where
    JPEG and resolution shortcuts are provably eliminated.
  \item Ablation studies on the effect of preprocessing normalisation
    on detector performance.
\end{itemize}

\paragraph{Uses to avoid.}
\begin{itemize}[leftmargin=1.5em, itemsep=2pt]
  \item Training detectors intended for deployment on images that have
    not been similarly preprocessed, without re-evaluating robustness
    to the omitted preprocessing.
  \item Using the \texttt{orig\_width} and \texttt{orig\_height} fields
    as training features; these encode source-level provenance and
    reintroduce a resolution-based shortcut.
  \item Any application requiring a larger or more demographically
    representative real-image corpus than COCO and ImageNet provide.
  \item Commercial deployment of detectors trained solely on this
    dataset without further validation, given the non-commercial
    licence restriction on the ImageNet component.
\end{itemize}

\paragraph{Impact of misuse.}
A detector trained under conditions that inflate performance
(e.g., on an uncorrected dataset) may provide false assurance
of reliability when deployed on images that do not share those
conditions.
We caution that \textbf{\datasetname} itself is a benchmark, not a
production detection system, and that all reported metrics should
be interpreted within the experimental conditions described in
Section~\ref{sec:experiments}.

%-------------------------------------------------------------------
\subsection*{A.6 \quad Distribution}
\label{ssec:ds_distribution}
%-------------------------------------------------------------------

\paragraph{Access.}
The dataset is publicly available at:\\
\texttt{https://huggingface.co/datasets/Shanmuk4622/ai-detection-dataset-v2}\\
All images, metadata, manifest, splits, validation report, and
generation configuration are accessible without authentication.

\paragraph{Licence.}
The dataset inherits the most restrictive licence applicable to any
of its components:
\begin{itemize}[leftmargin=1.5em, itemsep=2pt]
  \item \textbf{COCO images}: Flickr photographs under Creative Commons
    licences (primarily CC BY 2.0). Commercial use permitted.
  \item \textbf{ImageNet images}: ILSVRC terms of use.
    \emph{Non-commercial research use only.}
  \item \textbf{AI images}: Synthetic outputs of the six generators,
    each under its own open model licence. All six permit
    non-commercial research use.
\end{itemize}
The effective licence for the combined dataset is therefore
\textbf{non-commercial research use only}.
This is not legal advice; users are responsible for verifying
compliance with the licence terms of each component.

\paragraph{DOI and versioning.}
The dataset is tagged at version \texttt{v2.0} in the HuggingFace
repository.
A DOI will be assigned upon journal acceptance through the
HuggingFace dataset DOI programme.

\paragraph{Export controls.}
No export control or other legal restrictions beyond the licence
terms above are known to apply to this dataset.

%-------------------------------------------------------------------
\subsection*{A.7 \quad Maintenance}
\label{ssec:ds_maintenance}
%-------------------------------------------------------------------

\paragraph{Maintainer.}
The dataset is maintained by the first author,
\texttt{Shanmuk4622} on HuggingFace.
Contact: via the HuggingFace repository discussion tab.

\paragraph{Erratum policy.}
Errors in image data, metadata, or splits discovered after release
will be documented in the repository's \texttt{ERRATA.md} file and
will not silently overwrite existing files.
Corrected versions will be tagged with an incremented minor version
number (e.g., \texttt{v2.1}) with a changelog entry.

\paragraph{Updates.}
No scheduled updates are planned.
Extensions (additional generators, larger image counts, or
higher resolution variants) will be released as separate versioned
datasets rather than by modifying the released \texttt{v2.0} tag,
preserving the reproducibility of published results.

\paragraph{Reproducibility.}
The complete build is reproducible from:
\begin{enumerate}[leftmargin=1.5em, itemsep=2pt]
  \item The eleven generation notebooks (NB00--NB11) stored in the
    repository under \texttt{notebooks/}.
  \item The shared library \texttt{ai\_detect\_common.py}
    (pipeline version 1.2).
  \item The configuration file \texttt{config.json}, which records
    all hyperparameters including \texttt{MASTER\_SEED}~$= 42$.
\end{enumerate}
Changing \texttt{MASTER\_SEED} in NB00 and re-running the full
pipeline produces a statistically independent dataset with the same
structure, enabling sensitivity analysis over the random seed.

\paragraph{Legacy support.}
\texttt{v1.x} of the dataset (50,000 images, different repository)
remains available at its original HuggingFace URL for continuity
with any work that cited it prior to this paper.
It is not recommended for new work; \texttt{v2.0} supersedes it
in all respects.

% ─── New bibliography entry ─────────────────────────────────────
% Add to thebibliography block in related_work_and_references.tex

%\bibitem{gebru2021datasheets}
%T.~Gebru, J.~Morgenstern, B.~Vecchione, J.~W.~Vaughan, H.~Wallach,
%H.~Daum\'{e}~III, and K.~Crawford,
%``Datasheets for datasets,''
%\emph{Communications of the ACM}, vol.~64, no.~12, pp.~86--92, 2021.
%doi:10.1145/3458723.

%
%  New \bibitem key:
%    gebru2021datasheets
%
%  Required preamble additions:
%    \usepackage{enumitem}       % for \begin{enumerate}[label=\textbf{Q\arabic*.}]
%    \usepackage{xcolor}
%    \definecolor{shadecolor}{RGB}{245,245,245}
%    \usepackage{framed}         % for \begin{shaded}...\end{shaded} question boxes
%    (or replace shaded blocks with \noindent\textbf{...} if framed not available)
% ============================================================

\section{Dataset Documentation}
\label{sec:datasheet}

This appendix provides a structured documentation record for
\textbf{\datasetname} following the \emph{Datasheets for Datasets}
framework of Gebru~\etal~\cite{gebru2021datasheets}.
We address all seven canonical categories: motivation, composition,
collection process, preprocessing and labelling, intended uses,
distribution, and maintenance.
This documentation is also provided in machine-readable form as
\texttt{datasheet.json} in the dataset repository.

%-------------------------------------------------------------------
\subsection*{A.1 \quad Motivation}
\label{ssec:ds_motivation}
%-------------------------------------------------------------------

\paragraph{Purpose.}
\textbf{\datasetname} was created to support the development and
evaluation of AI-generated image detectors under conditions that
eliminate the two principal dataset biases identified in prior work:
JPEG compression asymmetry between real and AI images, and semantic
content mismatch between the two classes.
No comparable leak-free paired dataset existed at the time of
construction that covered more than four generators or more than one
backbone architecture family.

\paragraph{Creators.}
The dataset was constructed by Shanmukesh Bonala
(\texttt{Shanmuk4622}) as part of doctoral research at VIT-AP
University, under the supervision of Dr.~Deepthi Godavarthi.
The build pipeline was implemented across eleven version-controlled
Kaggle notebooks (NB00--NB11) available at the repository URL.

\paragraph{Funding.}
No external funding was received for this work.
Computational resources were provided by Kaggle's free-tier GPU
programme (dual NVIDIA Tesla T4).

%-------------------------------------------------------------------
\subsection*{A.2 \quad Composition}
\label{ssec:ds_composition}
%-------------------------------------------------------------------

\paragraph{What the dataset represents.}
\textbf{\datasetname} consists of 70,000 RGB images at $512 \times 512$
pixels, divided into a real class (10,000 photographs) and an AI
class (60,000 synthetic images produced by six text-to-image diffusion
models).
Each real image has exactly six AI partners, one per generator,
all generated from the same BLIP-2 caption derived from that real
image.

\paragraph{Instances.}
Each instance is a single image stored as a lossless PNG in an Apache
Parquet shard.
Metadata fields (listed in full in Table~\ref{tab:schema}) include a
globally unique image identifier, the pairing key
(\texttt{source\_real\_id}), label (0~=~real, 1~=~AI), generator key,
source dataset name, generation prompt, spatial dimensions, original
pre-crop dimensions, generation hyperparameters, caption model name,
pipeline version, SHA-256 digest of the PNG bytes, and a UTC creation
timestamp.

\paragraph{Split counts.}
The dataset is partitioned into a training split (49,392 images,
7,056 pairs), a validation split (10,122 images, 1,446 pairs), and a
test split (10,486 images, 1,498 pairs).
Splits are assigned deterministically at the pair level
(Section~\ref{ssec:splits}) so that no real image and its six AI
partners appear in different folds.

\paragraph{Label noise.}
All real images carry label~0 and all AI images carry label~1 by
construction.
No human annotation was involved in label assignment; labels are
determined entirely by the data source.
Label noise is therefore zero by design, but generation failures
(e.g., a model collapsing to a degenerate output) may produce
images that are visually uninformative.
One such group was identified during deduplication (four W\"{u}rstchen
images sharing identical byte content, flagged in the manifest via
shared SHA-256 digest) and is noted in Section~\ref{ssec:dedup}.

\paragraph{Sensitive content.}
Real images are sourced from MS-COCO and ImageNet validation splits.
COCO images are Flickr-sourced photographs under Creative Commons
licences; they may contain depictions of people, including children,
in public settings.
AI images are synthetic outputs and do not depict real individuals.
The dataset contains no sexually explicit content; generation
pipelines were run with \texttt{safety\_checker=None} solely to
prevent the safety filter from introducing a class-asymmetric
post-processing step, not to produce unsafe content.

\paragraph{Self-contained.}
Image bytes are embedded directly in the Parquet files; no external
URLs or file paths are required to reproduce the dataset.

%-------------------------------------------------------------------
\subsection*{A.3 \quad Collection Process}
\label{ssec:ds_collection}
%-------------------------------------------------------------------

\paragraph{Real images.}
5,000 images were sampled from the MS-COCO 2017 validation
split~\cite{lin2014coco} and 5,000 from the ImageNet-1k validation
split~\cite{deng2009imagenet}, downloaded via the canonical
HuggingFace dataset repositories.
No web scraping was performed; all real images were obtained from
their official public releases.
Sampling was deterministic: images were drawn in index order up to
the target count, with MASTER\_SEED~$= 42$ controlling any shuffled
ordering.

\paragraph{AI images.}
Synthetic images were generated programmatically during June 2026
on Kaggle GPU sessions (NVIDIA Tesla T4, 15.6~GB VRAM) using the
HuggingFace Diffusers library.
Six generation notebooks (NB03--NB08) ran in parallel across multiple
Kaggle accounts using a hash-based worker-partitioning scheme to
avoid shard collisions.
No human selection was applied to generated images; every output
produced by the pipeline was included, subject only to the
integrity checks described in NB10.

\paragraph{Captions.}
Captions were generated in NB02 by feeding each real image (after
canonical preprocessing) to BLIP-2 (\texttt{Salesforce/blip2-opt-2.7b})
running beam search (5 beams, max 60 new tokens, repetition penalty
1.3), then cleaning and token-capping the result as described in
Section~\ref{ssec:captions}.
No human captioning was performed.

\paragraph{Workers and consent.}
All data collection and generation was automated; no human workers
or crowdworkers were involved.

%-------------------------------------------------------------------
\subsection*{A.4 \quad Preprocessing, Cleaning, and Labelling}
\label{ssec:ds_preprocessing}
%-------------------------------------------------------------------

\paragraph{Preprocessing.}
Every image — real or AI — was passed through the canonical
preprocessing pipeline (Algorithm~\ref{alg:preprocess}) with no
branching on the class label.
The pipeline performs EXIF transpose, RGB conversion, centre-crop to
square, Lanczos resize to $512 \times 512$, JPEG equalisation at
quality factor 95 with 4:4:4 chroma subsampling, and lossless PNG
serialisation at compress level 6.

\paragraph{Quality checks.}
Every stored image was verified against the \texttt{png\_is\_canonical}
predicate (PNG format, RGB mode, $512 \times 512$ dimensions, no ICC
profile) and the stored SHA-256 digest.
Images failing these checks were excluded and the corresponding
generation run was re-executed.
Zero images in the released dataset fail either check.

\paragraph{Labelling.}
Labels are assigned automatically by provenance: images sourced from
COCO or ImageNet receive label~0; images output by any of the six
generators receive label~1.
No human labelling was performed.

\paragraph{Raw data availability.}
The raw (pre-preprocessed) real images are available from their
respective original dataset repositories.
Raw generator outputs prior to canonical preprocessing are not
retained in the repository; only the preprocessed canonical form
is stored.
The preprocessing code is fully open-source and deterministic, so
raw generator outputs can be reproduced from the stored prompts
and seeds.

%-------------------------------------------------------------------
\subsection*{A.5 \quad Uses}
\label{ssec:ds_uses}
%-------------------------------------------------------------------

\paragraph{Intended uses.}
\textbf{\datasetname} is intended for:
\begin{itemize}[leftmargin=1.5em, itemsep=2pt]
  \item Training and evaluating binary real-vs-AI image classifiers
    under bias-free conditions.
  \item Studying per-generator detectability gradients and
    cross-generator generalisation.
  \item Benchmarking novel detection architectures on a dataset where
    JPEG and resolution shortcuts are provably eliminated.
  \item Ablation studies on the effect of preprocessing normalisation
    on detector performance.
\end{itemize}

\paragraph{Uses to avoid.}
\begin{itemize}[leftmargin=1.5em, itemsep=2pt]
  \item Training detectors intended for deployment on images that have
    not been similarly preprocessed, without re-evaluating robustness
    to the omitted preprocessing.
  \item Using the \texttt{orig\_width} and \texttt{orig\_height} fields
    as training features; these encode source-level provenance and
    reintroduce a resolution-based shortcut.
  \item Any application requiring a larger or more demographically
    representative real-image corpus than COCO and ImageNet provide.
  \item Commercial deployment of detectors trained solely on this
    dataset without further validation, given the non-commercial
    licence restriction on the ImageNet component.
\end{itemize}

\paragraph{Impact of misuse.}
A detector trained under conditions that inflate performance
(e.g., on an uncorrected dataset) may provide false assurance
of reliability when deployed on images that do not share those
conditions.
We caution that \textbf{\datasetname} itself is a benchmark, not a
production detection system, and that all reported metrics should
be interpreted within the experimental conditions described in
Section~\ref{sec:experiments}.

%-------------------------------------------------------------------
\subsection*{A.6 \quad Distribution}
\label{ssec:ds_distribution}
%-------------------------------------------------------------------

\paragraph{Access.}
The dataset is publicly available at:\\
\texttt{https://huggingface.co/datasets/Shanmuk4622/ai-detection-dataset-v2}\\
All images, metadata, manifest, splits, validation report, and
generation configuration are accessible without authentication.

\paragraph{Licence.}
The dataset inherits the most restrictive licence applicable to any
of its components:
\begin{itemize}[leftmargin=1.5em, itemsep=2pt]
  \item \textbf{COCO images}: Flickr photographs under Creative Commons
    licences (primarily CC BY 2.0). Commercial use permitted.
  \item \textbf{ImageNet images}: ILSVRC terms of use.
    \emph{Non-commercial research use only.}
  \item \textbf{AI images}: Synthetic outputs of the six generators,
    each under its own open model licence. All six permit
    non-commercial research use.
\end{itemize}
The effective licence for the combined dataset is therefore
\textbf{non-commercial research use only}.
This is not legal advice; users are responsible for verifying
compliance with the licence terms of each component.

\paragraph{DOI and versioning.}
The dataset is tagged at version \texttt{v2.0} in the HuggingFace
repository.
A DOI will be assigned upon journal acceptance through the
HuggingFace dataset DOI programme.

\paragraph{Export controls.}
No export control or other legal restrictions beyond the licence
terms above are known to apply to this dataset.

%-------------------------------------------------------------------
\subsection*{A.7 \quad Maintenance}
\label{ssec:ds_maintenance}
%-------------------------------------------------------------------

\paragraph{Maintainer.}
The dataset is maintained by the first author,
\texttt{Shanmuk4622} on HuggingFace.
Contact: via the HuggingFace repository discussion tab.

\paragraph{Erratum policy.}
Errors in image data, metadata, or splits discovered after release
will be documented in the repository's \texttt{ERRATA.md} file and
will not silently overwrite existing files.
Corrected versions will be tagged with an incremented minor version
number (e.g., \texttt{v2.1}) with a changelog entry.

\paragraph{Updates.}
No scheduled updates are planned.
Extensions (additional generators, larger image counts, or
higher resolution variants) will be released as separate versioned
datasets rather than by modifying the released \texttt{v2.0} tag,
preserving the reproducibility of published results.

\paragraph{Reproducibility.}
The complete build is reproducible from:
\begin{enumerate}[leftmargin=1.5em, itemsep=2pt]
  \item The eleven generation notebooks (NB00--NB11) stored in the
    repository under \texttt{notebooks/}.
  \item The shared library \texttt{ai\_detect\_common.py}
    (pipeline version 1.2).
  \item The configuration file \texttt{config.json}, which records
    all hyperparameters including \texttt{MASTER\_SEED}~$= 42$.
\end{enumerate}
Changing \texttt{MASTER\_SEED} in NB00 and re-running the full
pipeline produces a statistically independent dataset with the same
structure, enabling sensitivity analysis over the random seed.

\paragraph{Legacy support.}
\texttt{v1.x} of the dataset (50,000 images, different repository)
remains available at its original HuggingFace URL for continuity
with any work that cited it prior to this paper.
It is not recommended for new work; \texttt{v2.0} supersedes it
in all respects.

% ─── New bibliography entry ─────────────────────────────────────
% Add to thebibliography block in related_work_and_references.tex

%\bibitem{gebru2021datasheets}
%T.~Gebru, J.~Morgenstern, B.~Vecchione, J.~W.~Vaughan, H.~Wallach,
%H.~Daum\'{e}~III, and K.~Crawford,
%``Datasheets for datasets,''
%\emph{Communications of the ACM}, vol.~64, no.~12, pp.~86--92, 2021.
%doi:10.1145/3458723.
